\section{Motivation and Problem}
\hspace{\parindent}
This dissertation is developed in the context of research on centralized multi-robot fleet
coordination conducted at Institute
for Systems and Computer Engineering, Technology and Science (INESC TEC), namely at Centre for Robotics in Industry and Intelligent Systems (CRIIS) and the Industry and Innovation Laboratory (iiLab). Within this context, a
centralized fleet coordination framework has been proposed to address multi-robot navigation under
realistic operational constraints, relying on a shared environment representation and a sequential
\emph{Time-Enhanced A*} (TEA*) planning component~\cite{matos2025criis,silver2005cooperative,moura2019temporal}.

In this framework, robots are planned sequentially according to a predefined priority order, which
ensures conflict avoidance by construction but makes the resulting solution highly sensitive to robot
ordering. Different priority permutations may lead to substantially different outcomes, particularly
in constrained environments with shared resources. As the number of possible orderings grows
factorially with fleet size, exhaustive exploration becomes computationally infeasible.

This limitation motivates the development of prioritization heuristics capable of generating
informative robot orderings without enumerating all possible permutations. By exploiting different
properties of the planning instance, these heuristics aim
to produce orderings that are more relevant than purely random priority selections.

However, no single heuristic is expected to perform best across all scenarios, since the usefulness of
a priority rule depends on the spatial configuration of the robots, the structure of the environment,
and the interactions between their paths. This motivates the investigation of a parallel planning
strategy in which multiple heuristic-generated orderings are evaluated within the same planning
cycle. By exploiting multithreading to execute several TEA* instances in parallel, the impact of
ordering sensitivity can be reduced while preserving compatibility with the existing TEA*
coordination framework.